\documentclass[11pt,a4paper]{article}
\usepackage[T1]{fontenc}
\usepackage{lmodern,amsmath,amssymb,amsthm,microtype,booktabs}
\usepackage[margin=25mm]{geometry}
\usepackage[colorlinks=true,urlcolor=blue,linkcolor=blue]{hyperref}
\newtheorem{theorem}{Theorem}
\newcommand{\dd}{\mathop{}\!\mathrm d}
\newcommand{\D}{\mathcal D}
\title{Cancellation zeros and the next recoil terms\\in the full Barnes amplitude}
\author{Alan Day \and Chengheng\thanks{AI research assistant (OpenAI).}}
\date{Working research note --- 30 October 2026}
\begin{document}\maketitle
\begin{abstract}
The two-branch Bessel leading term has infinitely many simple zeros
on the positive real parameter plane. We compute the next integrated
coefficients of the exact Barnes kernel and show that these auxiliary
zeros are lifted at order \(\nu^{-2}\). Thus a uniform relative error
dividing by the leading sum cannot include this set. Away from
cancellation, the second coefficient explains the previously observed
\(K^3/\nu^2\) scale. The integrated expansion is proved at fixed
parameters; a joint growing-\(K\) remainder remains open.
\end{abstract}

\section{Object and a fixed-parameter expansion}
Use radius one, \(u=v=0\), \(\omega=i\alpha\), \(K>0\),
real \(\alpha\), \(M=\nu^2/(2K)\). The object is the internal chiral
graph amplitude after the fixed-mass \(\kappa\to\infty\) limit.
No finite-\(G\) state, physical external insertion or normalized OTOC
is identified here. Set \(t=z+i(M+\nu)\),
\[
 b=\tfrac12-i(M-\nu),\quad c=\tfrac12-i(M+\nu),\quad
 C=e^{-\pi(M+\nu)}\frac{\Gamma(-i\nu/2)^4\Gamma(b)}{\Gamma(c)}.
\]
The normalized kernel, inherited from the preceding full-limit note, is
\[
 Q_\nu(t)=e^{-i\pi t}\frac{\Gamma(t-i\nu/2)^4}{\Gamma(-i\nu/2)^4}
 \frac{\Gamma(b-t)\Gamma(c)}{\Gamma(b)\Gamma(c+t)}
 \Gamma(\tfrac12-t+\alpha/2)\Gamma(\tfrac12-t-\alpha/2).
\]
Its contour runs upward from \(-N-i\infty\) to \(-N+i\eta\),
then right, with \(\eta>0\), integer \(N>|\alpha|/2+1/2\).
Write \(\mathcal I_\nu=(2\pi i)^{-1}\int Q_\nu\,\dd t=\mathfrak B/C\).
Define, with \(\D=(K/2)\partial_K\) at fixed \(\alpha\),
\begin{align*}
 J&=-K[e^{i\pi\alpha}K_\alpha(K)+i\pi I_\alpha(K)],\qquad
 S=K^2+\alpha^2-1,\\
 T&=\tfrac13[2S\D J+(K^2-\alpha^2+1)J],\qquad
 P=\tfrac16[9K^4-(\alpha^2-1)^2].
\end{align*}
\begin{theorem}
For fixed \(K,\alpha\), as \(\nu\to\infty\),
\[
 \mathcal I_\nu
 =e^{i(S/\nu+P/\nu^3)}\left(J+\frac{T}{\nu^2}\right)
  +O_{K,\alpha}(\nu^{-4}). \tag{1}
\]
The error is absolute, including at \(J=0\).
\end{theorem}

\newpage
\section{Derivation of the integrated coefficients}
Let
\[
 U(t)=e^{-2\pi it}(K/2)^{2t}
 \Gamma(\tfrac12-t+\alpha/2)\Gamma(\tfrac12-t-\alpha/2).
\]
On the logarithmically growing truncation \(R=\log\nu\),
Gamma-ratio expansions \cite{gamma} on safe sectors give
\[
 \log(Q_\nu/U)=\frac{L_1}{\nu}+\frac{L_2}{\nu^2}
 +\frac{L_3}{\nu^3}
 +O_{K,\alpha}((1+|t|)^5/\nu^4), \tag{2}
\]
where log ratios are continued from zero shift, and
\begin{align*}
 L_1&=4i(t^2-t),\\
 L_2&=\tfrac83 B_3(t)+4K^2t
     =\tfrac83t^3-4t^2+(\tfrac43+4K^2)t,\\
 L_3&=4iK^2t^2-\tfrac{8i}{3}(t^2-t)^2.
\end{align*}
For the four arguments \(-i\nu/2\), the Bernoulli terms through
\(B_4\) give the displayed cubic and quartic polynomials.
The leading \(b,c\) powers contribute
\((4K^2-2iK)t/\nu^2\); their quadratic shifts supply
\(2iKt/\nu^2+4iK^2t^2/\nu^3\). This explicitly cancels the imaginary
second-order term.

Exponentiating (2), the coefficient polynomials are
\[
 A_2=L_2+\tfrac12L_1^2
 =-8t^4+\tfrac{56}3t^3-12t^2+(\tfrac43+4K^2)t,
 \quad
 A_3=L_3+L_1L_2+\tfrac16L_1^3. \tag{3}
\]
Replace \(t^n\) under the auxiliary integral by \(\D^n\), keeping
the displayed coefficient powers of \(K\) to the left.
Since \(\D^2J=\D J+(S/4)J\), every result reduces to \(J,\D J\).
For a direct algebraic verification, if \(q=S/4\), then
\[
 \D(FJ+G\D J)=(\D F+qG)J+(F+\D G+G)\D J,\qquad
 \D q=K^2/4.
\]
This gives
\[
 A_2(\D)J=T-\tfrac12S^2J,\qquad
 A_3(\D)J-iS A_2(\D)J-\tfrac{i}{3}S^3J=iPJ. \tag{4}
\]
The third term after removal of the first phase is thus proportional
to \(J\), while the second contains its derivative.

For completeness, the preceding tail proof supplies, uniformly in
large \(\nu\) at fixed \(K,\alpha\),
\[
 |Q_\nu(-N-is)|\le C(1+s)^{2N}e^{-3\pi s},\quad
 \int_R^\infty|Q_\nu(x+i\eta)|\dd x\le6RH_R,
\]
with \(H_R=\sup_{[R,R+1]}|Q_\nu(x+i\eta)|\) once
\(R\ge\max(2,2K,2\max(\alpha^2/4-\eta^2,0))\).
At \(R=\log\nu\), (2) implies
\(H_R\le C\exp[-2R\log R+O_K(R)]\).
The vertical tail is \(o(\nu^{-4})\), the horizontal tail is smaller
than every fixed inverse power, and all polynomial moments of
\(|U|\) are integrable. The truncated exponentiation error is
bounded by \(C\nu^{-4}(1+|t|)^8|U|\).
Integration and (4) prove (1). The constants are not joint-\(K\) bounds.

\newpage
\section{The leading sum has simple real cancellation zeros}
It suffices to take \(\alpha\ge0\), since the complete \(J\) is even
in \(\alpha\). For \(K>0\), \(I_\alpha(K)\) and \(K_\alpha(K)\) are
positive \cite{limits}. Hence
\[
 \Re J=-K\cos(\pi\alpha)K_\alpha(K).
\]
A zero requires \(\cos(\pi\alpha)=0\). If \(\sin(\pi\alpha)=+1\),
the imaginary terms add. Therefore all positive-order zeros occur at
\[
 \alpha_n=2n+\tfrac32,\quad n=0,1,\ldots,\qquad
 K_{\alpha_n}(K_n)=\pi I_{\alpha_n}(K_n). \tag{5}
\]
There is exactly one \(K_n>0\) for each such \(\alpha_n\).
Indeed the ratio \(r(K)=K_\alpha(K)/I_\alpha(K)\) runs from infinity
to zero, and the standard Wronskian \cite{wronskian} gives
\[
 r'(K)=-\frac1{K I_\alpha(K)^2}<0.
\]
At the zero,
\[
 J'(K_n,\alpha_n)=-\frac{i}{I_{\alpha_n}(K_n)},\qquad
 T(K_n,\alpha_n)=-\frac{iK_nS_n}{3I_{\alpha_n}(K_n)}\ne0,
 \quad S_n=K_n^2+\alpha_n^2-1. \tag{6}
\]
Equation (1) then implies, at each fixed zero,
\[
 \mathcal I_\nu(K_n,\alpha_n)
 =-\frac{iK_nS_n}{3I_{\alpha_n}(K_n)\nu^2}
 e^{iS_n/\nu}+O_{K_n,\alpha_n}(\nu^{-4}). \tag{7}
\]
The exact integral therefore does not retain that auxiliary zero for
all sufficiently large \(\nu\). This is not a claim locating all
zeros of the finite-mass integral or excluding displaced zeros.

The zeros accumulate near the previously found magnitude crossover.
Let \(F(y)=\sqrt{1+y^2}-y\operatorname{arsinh}y\), with
positive root \(y_c\simeq1.50887956\).
Uniform Bessel asymptotics \cite{order} give, near \(y_c\),
\[
\log\frac{K_\alpha(K)}{\pi I_\alpha(K)}
=-2KF(\alpha/K)+O(K^{-1}).
\]
Monotonicity first traps the root there, and then
\[
 \alpha_n/K_n=y_c+O(\alpha_n^{-2}). \tag{8}
\]
These are auxiliary interference zeros, not physical critical points.

\begin{center}\begin{tabular}{rrr}\toprule
\(\alpha_n\)&\(K_n\) (approximate)&
\(\alpha_n^2(\alpha_n/K_n-y_c)\)\\\midrule
1.5&1.00000000&\(-0.0199790\)\\
3.5&2.32218535&\(-0.0205628\)\\
7.5&4.97178686&\(-0.0206769\)\\
15.5&10.27310967&\(-0.0207018\)\\
31.5&20.87670645&\(-0.0207076\)\\
63.5&42.08435037&\(-0.0207089\)\\\bottomrule
\end{tabular}\end{center}

\newpage
\section{Consequences for the joint recoil window}
Any relative error divided by \(J\) is undefined on (5), and (7)
shows that this is not just a harmless shared zero. A prospective
uniform result should instead control the complete complex amplitude,
for example against the positive branch envelope
\[
 \mathcal E(K,\alpha)=K[K_\alpha(K)+\pi I_\alpha(K)].
 \tag{9}
\]
A bound relative to \(\mathcal E\) has not yet been proved for the
joint limit. It would preserve absolute control near interference
and imply a relative estimate only where \(|J|\) is bounded below
by a suitable fraction of \(\mathcal E\).

Away from cancellation, set
\[
 c_2(K,\alpha)=T(K,\alpha)/J(K,\alpha)-T(K,0)/J(K,0).
\]
For \(\alpha=Ky\), fixed \(y>0\) separated from \(y_c\), Bessel
derivative asymptotics at \emph{fixed order when differentiating}
give the coefficient limits
\[
 \frac{c_2(K,Ky)}{K^3}\longrightarrow
 \begin{cases}
 ((1+y^2)^{3/2}-1)/3,&0<y<y_c,\\
 -((1+y^2)^{3/2}+1)/3,&y>y_c.
 \end{cases} \tag{10}
\]
The reference \(y=0\) contributes \(+1/3\).
Equation (10) is an asymptotic of an explicitly known coefficient,
not of the unbounded remainder in (1).
It identifies the candidate small scale
\[
 K^3/\nu^2=\nu^4/(8M^3).
\]
Thus \(M\gg\nu^{4/3}\) suppresses this coefficient, while
\(M\ll\nu^2\) ensures \(K\gg1\).
This agrees with the earlier candidate recoil window but does not
prove its sufficiency. The first phase \(K^2/\nu\) can be large in
that window, so a joint analysis must retain it rather than truncate it.

At fixed parameters, since \(J(K,0)\ne0\), (1) also yields the
absolute normalized-ratio approximation
\[
 \frac{\mathfrak B(K,\alpha)}{\mathfrak B(K,0)}
 =
 e^{i[\alpha^2/\nu-(\alpha^4-2\alpha^2)/(6\nu^3)]}
 \frac{J(K,\alpha)+T(K,\alpha)/\nu^2}
 {J(K,0)+T(K,0)/\nu^2}+O_{K,\alpha}(\nu^{-4}). \tag{11}
\]
This form retains the derivative term even when the numerator's
leading \(J\) vanishes.

Reusing the frozen full-integral data at \(\nu=4K^3\) gives:
\begin{center}\begin{tabular}{rrrrr}\toprule
\(K\)&\(y\)&first phase only&plus \(c_2/\nu^2\)&formula (11)\\\midrule
4&2&\(3.89\,10^{-3}\)&\(4.02\,10^{-5}\)&\(6.35\,10^{-6}\)\\
8&1&\(7.43\,10^{-5}\)&\(7.71\,10^{-8}\)&\(6.64\,10^{-9}\)\\
8&2&\(4.93\,10^{-4}\)&\(1.27\,10^{-6}\)&\(1.02\,10^{-7}\)\\
16&2&\(6.19\,10^{-5}\)&\(3.97\,10^{-8}\)&\(1.60\,10^{-9}\)\\\bottomrule
\end{tabular}\end{center}
Entries are relative numerical differences; no coefficient was fitted.

\newpage
\section{A joint positive-tail ingredient at an \(O(K)\) cutoff}
The earlier recurrence bound began at \(A=O(\alpha^2)\). This
threshold can be bridged under the same candidate hierarchy,
without raising the contour height. Fix \(Y\ge0\), \(K\ge1\),
\(|\alpha|\le KY\), and \(\eta=1\). Set
\[
 D_Y=\max(2,Y^2/2),\quad c_Y=8+2Y,\quad
 A=\max(2,2K,2d),\quad d=\max(\alpha^2/4-1,0).
\]
Assume
\[
 \nu^2\ge32(D_Y+1)K^3,\qquad R\ge c_YK. \tag{12}
\]
The exact recurrence from the preceding note is
\[
 \mathcal R_\nu(t)=\frac{Q_\nu(t+1)}{Q_\nu(t)}
 =\frac{(t-i\nu/2)^4}
 {[(t+\tfrac12-i\nu)^2+M^2][(t+\tfrac12)^2-\alpha^2/4]}.
\]
If \(R\le A+1\), write \(a=\nu/2-1\), \(L=M-\nu+1\).
The earlier denominator bounds, \(L\ge\nu^2/(4K)\), and
\(x^2-d\ge x^2/2\) give, for \(R\le x\le A+1\),
\begin{align*}
 |\mathcal R_\nu(x+i)|
 &\le \frac{2(x^2+a^2)^2}{L^2x^2}\\
 &\le2\left(\frac{A+1}{L}+\frac{a^2}{RL}\right)^2
 \le2\left(\frac18+\frac18\right)^2=\frac18 .
\end{align*}
Here \(A+1\le(D_Y+1)K^2\), so the first fraction is at most
\(4(D_Y+1)K^3/\nu^2\); the second is at most \(K/R\).
Assumption (12) also implies the earlier condition \(\nu\ge8K\).
For \(x\ge A\), the already proved step bound
\(|\mathcal R_\nu(x+i)|\le[x/(x+1)]^{3/2}\) applies.

Thus the geometrically contracting transition reaches the
polynomially bounded far tail. If
\(n_0=\max(0,\lceil A-R\rceil)\), the first \(n_0\) unit steps
cost at most \(8^{-n_0}\), and the remaining tail is bounded by
\(6(R+n_0)H_R8^{-n_0}\). Summing the finite transition gives the
convenient universal estimate
\[
 \int_R^\infty |Q_\nu(x+i)|\,\dd x
 \le10(R+1)H_R,\qquad
 H_R=\sup_{R\le x\le R+1}|Q_\nu(x+i)| . \tag{13}
\]
If \(R\ge A\), the original far-tail bound immediately implies (13).

This is a joint-\(K\) positive-tail lemma: the base cutoff need only
be \(O_Y(K)\), rather than \(O_Y(K^2)\).
It does not bound \(H_R\) relative to the branch envelope, nor the
negative contour with \(N=O(K)\). Those two tasks, together with
the integrated local remainder after phase retention, remain
necessary for a full joint asymptotic theorem.
Eighteen finite samples check the new contraction inequality,
including masses at the threshold in (12).

\newpage
\section{Direct checks at zeros and next proof requirements}
At the first two auxiliary zeros we directly integrate the exact
\(Q_\nu\), retaining both complex contour segments.
For the leading nonzero prediction in (7), we find
\begin{center}\begin{tabular}{rrr}\toprule
\(\alpha\)&\(\nu\)&relative error against (7)'s leading term\\\midrule
1.5&64&\(3.97013\,10^{-4}\)\\
1.5&128&\(9.91999\,10^{-5}\)\\
1.5&256&\(2.47966\,10^{-5}\)\\
3.5&64&\(9.38078\,10^{-4}\)\\
3.5&128&\(2.32015\,10^{-4}\)\\
3.5&256&\(5.78474\,10^{-5}\)\\\bottomrule
\end{tabular}\end{center}
The segment cancellation ratio reaches about \(1.70\,10^7\).
We use 20/32 Gauss nodes per segment and 60/80 digits, with a
100-digit comparison at the largest \(\nu\).
The latter also changes \((X,T,\eta)=(36,18,1)\) to
\((44,24,0.75)\). Refinements compare the complete complex sum.
These are empirical diagnostics, not interval certification.

All 65 checks pass: 3 symbolic operator identities, 4 compact
remainder-decrease checks, 18 zero/Wronskian/lifting identities,
6 quadrature refinements, 4 precision/contour comparisons,
4 zero-lifting error decreases, 8 improvements in reused joint
samples, and 18 joint-tail samples.
Six compact residuals and six reused fixed-parameter
second-order residuals are also stored.
Code, versions, parameters, complex segments, and input SHA256
hashes are in the accompanying project records.

The next uniform theorem must use an absolute or branch-envelope
error and keep the derivative correction through cancellation.
A promising route is to split the two analytic contributions before
bounding their sum, then prove contour estimates with explicit
\(K,\alpha\) dependence after retaining the first phase.
The positive-tail transition is now controlled by (13); the base
size and the negative contour still require estimates with explicit
parameter dependence.
The joint remainder, transverse variables, physical external energy
cycles and normalization, disconnected baseline, finite-\(G\) limit,
and microscopic dictionary all remain open.
This is an author self-reviewed proof draft.

\begin{thebibliography}{9}
\bibitem{gamma} NIST DLMF, \S5.11, Gamma asymptotics,
\url{https://dlmf.nist.gov/5.11}.
\bibitem{wronskian} NIST DLMF, 10.28.2, Bessel Wronskian,
\url{https://dlmf.nist.gov/10.28}.
\bibitem{limits} NIST DLMF, \S10.30, \S10.32, and \S10.40,
positivity and limiting forms,
\url{https://dlmf.nist.gov/10.32}.
\bibitem{order} NIST DLMF, \S10.41(ii), including derivative expansions,
\url{https://dlmf.nist.gov/10.41}.
\end{thebibliography}
\end{document}
